\section{Methodology}~\label{sec:methodology}

\subsection{Overview}~\label{subsec:method-overview}

This section sets out the research methodology: the aims and questions the work
addresses, the method used to answer each research question, the schedule and
milestones against which progress is tracked, the ethical obligations that apply,
and the threats to the validity of the results.

The research is conducted as \textbf{two sequential experimental studies} that
share one design and one artefact.

\begin{itemize}
    \item \textbf{Study 1 (RQ1)} validated the monitoring pipeline architecture
    using replayed telemetry from a publicly available, pre-labelled drone
    dataset. Because ground-truth health state was supplied by the dataset, the
    only thing under test was whether the pipeline could collect, store and
    surface health state in real time.
    \item \textbf{Study 2 (RQ2)} removes the labels. The same pipeline is
    connected to a physical DJI Tello drone, fault thresholds are derived
    empirically from the drone's own baseline telemetry, and Prometheus alerting
    rules must determine health state autonomously from raw sensor readings.
\end{itemize}

There is a reason for running the studies in this order, and it is not only that
the project spans two trimesters. Study 1 eliminates the pipeline architecture as
a source of error, so any detection failure in Study 2 can be attributed to the
detection logic, meaning the thresholds and alerting rules, rather than to the
ingestion, storage or query layers beneath them. Establishing the instrument
before testing the hypothesis is standard practice in experimental software
engineering \citep{wohlin2012}, and it is what allows Study 2 to isolate
threshold calibration as the variable of interest.

\subsubsection{Candidate Methods and Selection}

Five methods were considered. Table~\ref{tab:methods} compares them and states
why each was or was not adopted.

\begin{table}[htbp]
\centering
\caption{Comparison of Candidate Research Methods}
\label{tab:methods}
\footnotesize
\begin{tabular}{|p{2.4cm}|p{3.4cm}|p{3.4cm}|p{4.6cm}|}
\hline
\textbf{Method} & \textbf{Strengths} & \textbf{Weaknesses} & \textbf{Decision and reason} \\
\hline
Experimental Research \citep{wohlin2012} & Produces real, measurable evidence; repeatable under controlled conditions & May not generalise beyond the test environment; needs controllable fault conditions & \textbf{Selected (primary).} The research questions ask for latency and false positive figures, which only controlled measurement on a running system can produce \citep{yousuf2024, paradeshi2022}. \\
\hline
Design Science Research \citep{hevner2004} & Produces a reusable artefact; iterative build and evaluate cycle & Evaluation can be subjective unless criteria are fixed in advance & \textbf{Selected (secondary).} The monitoring pipeline is a contribution in itself, not only an instrument, so it needs a method that evaluates artefacts \citep{boncea2016}. \\
\hline
Simulation / Modelling & Safe and repeatable; no hardware cost & Does not reproduce real hardware behaviour, Wi-Fi variability, thermal drift or battery ageing & \textbf{Rejected.} The gap concerns the absence of evidence from physical mobile platforms. Simulating the platform would reproduce the gap instead of closing it. \\
\hline
Case Study / Field Observation & High realism; captures operational context & No controlled fault injection; poor repeatability & \textbf{Rejected.} RQ2.3 needs known fault onset times, which uncontrolled observation cannot supply. \\
\hline
Literature Review & Establishes grounding; identifies gaps and prior approaches & Produces no original data; cannot validate an artefact & \textbf{Supporting input only.} Used to establish the gap and select tooling and metrics. \\
\hline
\end{tabular}
\end{table}

\noindent\textbf{Experimental research (primary).} The primary method is a
controlled, repeated-measures fault-injection experiment. Baseline flights under
normal conditions are the control condition and flights with deliberately induced
degradation are the treatment conditions. The independent variable is the class
of induced fault. The dependent variables are detection outcome, detection
latency, and alerts fired when no fault was induced. A latency figure needs a
known fault onset time. A false positive rate needs a period of flight certified
free of faults. Both are properties of a controlled experiment and neither
survives in uncontrolled observation.

\noindent\textbf{Design science research (secondary).} The monitoring pipeline is
itself a contribution, and DSR provides an established structure for building and
evaluating such artefacts \citep{hevner2004}. \citet{boncea2016} used the same
approach to propose the Prometheus-based architecture this thesis builds on. The
artefact is the pipeline: telemetry source $\rightarrow$ Python exporter
$\rightarrow$ Prometheus $\rightarrow$ alerting rules $\rightarrow$ Grafana
dashboard. It was developed across two DSR iterations, one per study.

\subsubsection{Experimental Platform and Tooling}

\begin{table}[htbp]
\centering
\caption{Required Resources}
\label{tab:resources}
\small
\begin{tabular}{|>{\raggedright\arraybackslash}p{3.5cm}|p{6cm}|p{4.3cm}|}
\hline
\textbf{Resource} & \textbf{Purpose} & \textbf{Access / version} \\
\hline
DJI Tello drone & Representative RAS for live telemetry collection (Study 2) & Provided by Deakin University \\
\hline
Laptop / PC & Runs exporter, Prometheus, Grafana and session logger on one host & Personally owned; Windows 11 \\
\hline
Wi-Fi (2.4\,GHz) & Tello SDK communication over UDP & Direct Tello access point \\
\hline
Prometheus & Time-series collection, storage and alert rule evaluation & Open-source \\
\hline
Grafana & Dashboard visualisation of telemetry and alert state & Open-source \\
\hline
Python + \texttt{djitellopy} & Telemetry extraction and Prometheus exporter & Open-source; Python 3.11.9 \\
\hline
Drone Telemetry Tampering Dataset v2 & Pre-labelled telemetry for Study 1 & Public, via Kaggle \citep{ekanayakadevlk2024}; CC BY-SA 4.0 \\
\hline
\end{tabular}
\end{table}

Prometheus was selected for its pull-based scrape model, native time-series
storage and PromQL query language, which supports the threshold-based alerting
RQ2 requires. It is also the framework most widely reported in comparable
monitoring architectures \citep{boncea2016}, which makes these results directly
comparable to that prior work. Grafana was selected as the visualisation layer
because it integrates natively with Prometheus and is likewise established in the
reviewed literature \citep{benhafaiedh2022}. The exporter is written in Python
for its suitability to telemetry handling and its access to \texttt{djitellopy}.

A push-based stack (Telegraf/InfluxDB) would avoid the scrape-interval latency
floor discussed below, and log-oriented tooling (ELK) suits event data better
than numeric time series. Both were rejected, the former for being less widely
evidenced in the RAS monitoring literature and the latter for being poorly
matched to continuously sampled numeric telemetry. Managed cloud IoT platforms
were rejected on principle, since the gap this thesis addresses is the absence of
monitoring approaches that work \textit{without} cloud dependency.

One trade-off is worth stating early. Prometheus' pull model imposes a hard lower
bound on detection latency equal to the scrape interval, however quickly the
exporter updates. Study 1 showed this directly: metric values updated at the
exporter every second, but observed detection latency was governed by the
15-second scrape cycle. Study 2 therefore uses a \SI{1}{\second} scrape interval,
accepting the higher storage and CPU cost in exchange for a latency floor
compatible with the targets in Section~\ref{subsec:method-rq2}.

\subsection{Aims}~\label{subsec:method-aims}

The research has four aims, which together define what the two studies are
intended to produce.

\begin{enumerate}
    \item \textbf{Design a monitoring pipeline for remote autonomous systems}
    that collects, stores and surfaces health telemetry in real time using
    open-source components and no cloud dependency.
    \item \textbf{Establish which telemetry metrics represent the operational
    health} of a remote autonomous system, and demonstrate that the pipeline can
    track transitions between healthy and faulted states.
    \item \textbf{Derive fault detection thresholds empirically} from a
    platform's own measured baseline behaviour, rather than assuming values
    carried over from other domains, and state the derivation procedure so that
    it transfers to other platforms.
    \item \textbf{Evaluate the resulting detection performance} on physical
    hardware under controlled fault conditions, reporting detection rate,
    detection latency and false positive rate against targets fixed in advance.
\end{enumerate}

Aims 1 and 2 are addressed by Study 1 and answer RQ1. Aims 3 and 4 are addressed
by Study 2 and answer RQ2.

\subsection{Research Questions}~\label{subsec:method-rqs}

Two research questions guide the study, each decomposed into measurable
sub-questions.

\vspace{0.5em}
\noindent\textbf{RQ1:} \textit{How can monitoring frameworks be used to assess
and maintain the health of remote autonomous systems?}

\begin{itemize}
    \item \textbf{RQ1.1:} Which telemetry metrics must a monitoring pipeline
    expose in order to represent the operational health state of a remote
    autonomous system?
    \item \textbf{RQ1.2:} Can a pull-based time-series monitoring framework track
    transitions between healthy and faulted states in RAS telemetry in real time?
\end{itemize}

\vspace{0.5em}
\noindent\textbf{RQ2:} \textit{How can system metrics such as battery level,
connectivity and performance be used to detect faults in remote autonomous
systems?}

\begin{itemize}
    \item \textbf{RQ2.1:} What metric thresholds reliably distinguish normal
    operation from fault conditions in a physical drone system, and how should
    those thresholds be derived?
    \item \textbf{RQ2.2:} How accurately can Prometheus alerting rules detect
    fault conditions from raw telemetry without pre-labelled data?
    \item \textbf{RQ2.3:} What are the detection latency and false positive rate
    of the threshold-based alerting system under controlled fault scenarios?
\end{itemize}

These questions target the gap established in Section~\ref{subsec:lit-gap}. The
literature reviewed rarely integrates real-time telemetry pipelines with physical
autonomous systems operating outside controlled industrial settings
\citep{boncea2016, devi2024}, and where threshold-based detection is used,
thresholds are typically assumed rather than measured from the platform under
observation \citep{yousuf2024}. RQ1 addresses the integration half of that gap.
RQ2 addresses the calibration half, which is the more consequential of the two.
Preliminary testing produced a concrete instance of why: an assumed
\SI{60}{\celsius} temperature threshold, carried over from industrial motor
monitoring conventions, produced a sustained false positive against real Tello
telemetry, because the Tello's normal operating range differs from the
assumption. That observation is the empirical motivation for the
baseline-calibration procedure in Section~\ref{subsec:method-rq2}.

Table~\ref{tab:design-matrix} maps each sub-question to the method, data,
instrument, analysis and success criterion that address it.

\begin{table}[htbp]
\centering
\caption{Research Design Matrix}
\label{tab:design-matrix}
\footnotesize
\begin{tabular}{|p{1.4cm}|p{2.6cm}|p{2.6cm}|p{3.1cm}|p{3.4cm}|}
\hline
\textbf{Sub-\newline question} & \textbf{Data source} & \textbf{Instrument} & \textbf{Analysis} & \textbf{Success criterion} \\
\hline
RQ1.1 & Drone Telemetry Tampering Dataset v2 & Python exporter metric schema & Mapping of exposed metrics to RAS health dimensions & All identified health dimensions represented by at least one exposed metric \\
\hline
RQ1.2 & Dataset v2 (six labelled anomaly phases) & Exporter $\rightarrow$ Prometheus $\rightarrow$ Grafana & PromQL transition counting against dataset labels & 100\% of labelled fault-phase transitions reflected in \texttt{ras\_health\_status} \\
\hline
RQ2.1 & 5 Tello baseline flight sessions & Live Tello SDK exporter at \SI{1}{\second} scrape & Per-metric distribution analysis; percentile and operational-limit derivation & Threshold defined for every monitored metric, each traceable to baseline data \\
\hline
RQ2.2 & 9 Tello fault-injection runs (3 types $\times$ 3 reps) & Prometheus alerting rules (\texttt{rules.yml}) & Confusion matrix; detection rate and precision & $\geq$8 of 9 induced faults detected by the correct rule \\
\hline
RQ2.3 & Same 9 runs plus 5 baseline sessions & Alert firing times vs.\ raw telemetry event log & Latency distribution (median, IQR, max); false alerts per unit flight time & Median latency $\leq$\SI{5}{\second}, max $\leq$\SI{10}{\second}; $\leq$1 false alert per \SI{10}{\minute} normal flight \\
\hline
\end{tabular}
\end{table}

\subsection{RQ1 Methodology}~\label{subsec:method-rq1}

Study 1 tests whether the pipeline can \textit{observe} health state. It does not
test whether the pipeline can \textit{determine} health state, which is RQ2's
concern.

\noindent\textbf{Data source.} Telemetry came from the Drone Telemetry Tampering
Dataset v2 \citep{ekanayakadevlk2024}, a publicly available corpus of drone
flight telemetry containing six labelled anomaly phases interleaved with normal
operation. A pre-labelled dataset was chosen deliberately: because each row
carries a ground-truth health state, any discrepancy between the dataset labels
and what the pipeline surfaced is attributable to the pipeline and nothing else.

\noindent\textbf{Instrument.} A Python exporter reads the dataset row by row and
republishes each record as Prometheus metrics on an HTTP endpoint, replaying the
recorded flight as though it were arriving live. Alongside the raw telemetry the
exporter publishes \texttt{ras\_health\_status}, a composite state metric taken
directly from the dataset label. Prometheus scrapes the endpoint on a fixed
interval and Grafana renders the metrics and the health state.

\noindent\textbf{Protocol.} The full dataset is replayed end to end in a single
continuous run at a \SI{15}{\second} scrape interval. Metric coverage is checked
throughout with \texttt{up} and per-metric \texttt{absent()} queries. Every
labelled phase transition in the source data is then compared against the
corresponding change in \texttt{ras\_health\_status} as recorded in Prometheus.

\noindent\textbf{Analysis.} RQ1.1 is answered by mapping each exposed metric to
the health dimensions identified in the literature review, and checking that
every dimension is represented. RQ1.2 is answered by counting labelled phase
transitions in the dataset and the corresponding transitions recorded in
Prometheus; a transition is counted as tracked only if it appears in
\texttt{ras\_health\_status} within one scrape interval of the labelled change.

\begin{table}[htbp]
\centering
\caption{Study 1 Evaluation Criteria (RQ1)}
\label{tab:eval-rq1}
\small
\begin{tabular}{|p{2.8cm}|p{5.4cm}|p{2.6cm}|p{3.0cm}|}
\hline
\textbf{Criterion} & \textbf{Operational definition} & \textbf{Target} & \textbf{Measured by} \\
\hline
Metric coverage & All defined health metrics present and non-stale in Prometheus for the full replay & 100\% of metrics & \texttt{up}, \texttt{absent()} \\
\hline
Transition fidelity & Labelled fault-phase transitions reflected by a corresponding change in \texttt{ras\_health\_status} & 100\% of transitions & PromQL transition count vs.\ dataset labels \\
\hline
Detection latency & Delay between a metric value changing at the exporter and that change being queryable in Prometheus & $\leq$ 1 scrape interval & Exporter log vs.\ Prometheus sample timestamps \\
\hline
Dashboard completeness & All monitored metrics and the health-state panel visible on one screen without scrolling & All metrics plotted & Dashboard inspection \\
\hline
\end{tabular}
\end{table}

\noindent\textbf{Limitation carried into Study 2.} Because health state was read
from a dataset label rather than computed from raw values, Study 1 says nothing
about whether the system can recognise a fault it has not been told about. That
limitation is what Study 2 exists to address, and it is the reason the two
studies are sequential.

\subsection{RQ2 Methodology}~\label{subsec:method-rq2}

Study 2 removes the labels. The same pipeline is connected to a physical DJI
Tello, and health state must be computed by Prometheus alerting rules from raw
telemetry alone.

\subsubsection{Platform Telemetry Constraints}

The Tello's available telemetry differs materially from that of the GPS-equipped
platforms described in much of the reviewed literature, and from the fields in
the Study 1 dataset. The standard Tello carries \textbf{no GNSS receiver}, so
geospatial metrics are unavailable, and Wi-Fi link quality is exposed as a
signal-to-noise ratio query rather than a continuously sampled RSSI value.
Table~\ref{tab:tello-fields} lists the fields the SDK does expose. This
constraint shapes the metric set and returns as a threat to external validity in
Section~\ref{subsec:method-threats}.

\begin{table}[htbp]
\centering
\caption{Tello SDK Telemetry Fields Used in Study 2}
\label{tab:tello-fields}
\small
\begin{tabular}{|p{3.6cm}|p{4.4cm}|p{1.6cm}|p{4.2cm}|}
\hline
\textbf{Field} & \textbf{Accessor} & \textbf{Unit} & \textbf{Role} \\
\hline
Battery level & \texttt{get\_battery()} & \% & Primary depletion fault signal \\
\hline
Flight time & \texttt{get\_flight\_time()} & s & Session duration; depletion rate denominator \\
\hline
Height & \texttt{get\_height()} & cm & Barometrically derived height; drift indicator rather than measured altitude \\
\hline
ToF distance & \texttt{get\_distance\_tof()} & cm & Primary altitude-hold signal \\
\hline
Barometer & \texttt{get\_barometer()} & cm & Altitude drift detection \\
\hline
Temperature (high / low) & \texttt{get\_highest\_
temperature()}, \texttt{get\_lowest\_
temperature()} & \si{\celsius} & Thermal state; settles toward an operating point in flight \\
\hline
Acceleration $x,y,z$ & \texttt{get\_acceleration
\_x/y/z()} & 0.001g & Derived instability metric \\
\hline
Attitude (pitch, roll, yaw) & \texttt{get\_pitch()}, \texttt{get\_roll()}, \texttt{get\_yaw()} & deg & Derived instability metric \\
\hline
Velocity $x,y,z$ & \texttt{get\_speed\_x/y/z()} & dm/s & Recorded for completeness; near zero at \SI{1}{\hertz} even during commanded translation \\
\hline
Wi-Fi SNR & \texttt{query\_wifi\_signal
\_noise\_ratio()} & dB & Link degradation fault signal \\
\hline
\end{tabular}
\end{table}

\subsubsection{Data Collection Protocol}

All sessions use a \SI{1}{\second} Prometheus scrape interval. The exporter,
Prometheus and session logger run on a single host, so every timestamp comes from
one system clock and clock synchronisation between devices cannot introduce
measurement error.

\noindent\textbf{Phase 1: Baseline (RQ2.1).} Five flight sessions under normal
operating conditions, each of four minutes. Flights follow a fixed manoeuvre
script (take-off, hover, translation, rotation, hover, land) so sessions are
comparable. No faults are induced and no interventions are made. The purpose is
to establish the empirical distribution of every monitored metric under normal
operation, which is what the thresholds are derived from.

Three controls hold the sessions comparable to one another. Each session begins
on a fully charged battery, because percentage-per-minute depletion is not
linear across the discharge curve and a part-charged start samples a different
region of it. Batteries are alternated between sessions so that battery identity
is not confounded with session order. The aircraft is returned to ambient
temperature before each session, since residual heat from a preceding flight
displaces the thermal distribution of the next one.

Four-minute sessions sample the upper region of the discharge curve only. The
lower region, approaching the automatic low-battery landing, is characterised by
the Phase 2 depletion runs, which continue until that failsafe engages.

\noindent\textbf{Phase 2: Fault injection (RQ2.2, RQ2.3).} Nine runs, made up of
three fault classes with three repetitions each.

\begin{enumerate}
    \item \textbf{Battery depletion.} Flight continues until the Tello's
    automatic low-battery landing engages. Tests whether the depletion rule fires
    with enough lead time before the failsafe.
    \item \textbf{Link degradation.} Separation between drone and controller is
    increased progressively and line of sight is deliberately obstructed,
    degrading Wi-Fi SNR. Tests the connectivity rule.
    \item \textbf{Thermal and load stress.} Consecutive sorties flown without a
    cool-down interval, on a part-discharged battery, with continuous high-speed
    manoeuvring and no hover rest. Baseline measurement established that onboard
    temperature settles toward an operating point rather than rising under load,
    so aggressive manoeuvring alone does not produce a thermal fault, whereas
    repeated sorties do. Tests the performance-degradation rules.
\end{enumerate}

\noindent\textbf{Ground-truth marking.} Live telemetry carries no labels, so
fault onset has to be established independently of the alerting system being
measured. Two timestamps are recorded per run: $t_{\text{induce}}$, the moment
the operator begins inducing the fault, written to the session log as an event
marker by the control script; and $t_{\text{onset}}$, the first crossing of the
defined threshold in the raw \SI{1}{\hertz} telemetry log that persists for at
least the hold duration applied by the rule, worked out retrospectively during
analysis.

The hold duration is part of the definition because a shorter crossing cannot
raise an alert, and anchoring latency to one would measure the interval between
two unrelated events. Crossings below that duration are recorded separately as
transient crossings, since their frequency indicates how much of the false
positive suppression is attributable to the hold rather than to the threshold
value.

Detection latency is measured from $t_{\text{onset}}$. Measuring from
$t_{\text{induce}}$ would fold the drone's physical response time and the
operator's reaction into the figure, and neither of those is part of what RQ2.3
asks about. $t_{\text{induce}}$ is still recorded, as context and to confirm each
run's fault was induced as intended.

\noindent\textbf{Data recorded per session.} Each session is logged under four
headings. The first three make a session identifiable and comparable to the
others; the fourth is the measurement itself.

\begin{itemize}
    \item \textbf{Session identity.} A unique session identifier, and the date
    and time of the flight.
    \item \textbf{Conditions.} The venue, whether the flight was indoors or
    outdoors, and the ambient temperature.
    \item \textbf{Hardware state.} The Tello firmware version, the battery
    identifier, and that battery's approximate cycle count.
    \item \textbf{Measurements.} The full \SI{1}{\hertz} raw telemetry log,
    every Prometheus alert firing and resolution timestamp, the operator event
    markers, and the Prometheus system logs.
\end{itemize}

Conditions and hardware state are recorded because they vary between sessions
and affect the readings. Ambient temperature shifts the thermal baseline, and
battery age changes the depletion rate. Logging both means an outlier session
can be explained rather than discarded.

\subsubsection{Threshold Derivation Procedure}~\label{subsec:method-thresholds}

Study 1's main limitation was that health state came from dataset labels instead
of being computed. Removing that dependency makes threshold derivation the
central methodological problem of this thesis, not an implementation detail.

A single uniform rule will not work, because the monitored metrics do not behave
the same way as one another. Battery level, for example, falls steadily
throughout every flight by design, so a percentile of its distribution would say
nothing useful about when the drone is actually at risk.
Table~\ref{tab:threshold-rules} groups metrics by how they behave and gives each
group its own derivation rule and rationale.

\begin{table}[htbp]
\centering
\caption{Threshold Derivation Rules by Metric Class}
\label{tab:threshold-rules}
\footnotesize
\begin{tabular}{|p{2.4cm}|p{3.1cm}|p{4.2cm}|p{4.1cm}|}
\hline
\textbf{Metric class} & \textbf{Metrics} & \textbf{Derivation rule} & \textbf{Rationale} \\
\hline
Monotonic / depleting & Battery level & Operational limit plus safety margin: alert above the firmware auto-land point, with the margin sized from the mean depletion rate observed in Phase 1 & The quantity declines by design, so no percentile of its distribution is meaningful. What matters is keeping enough lead time to act before the failsafe. \\
\hline
Stationary / bounded & Wi-Fi SNR & Empirical 1st or 99th percentile of the pooled Phase 1 distribution, per metric direction & Percentiles assume nothing about normality. Telemetry from a mobile platform is skewed and multi-modal across flight modes, so mean $\pm k\sigma$ would put the bound in the wrong place. \\
\hline
Settling & High and low onboard temperature & 99th percentile of the settled region, taken as the final \SI{90}{\second} of each Phase 1 session, per metric direction & The quantity converges toward an operating point rather than varying about a mean, so most samples in a short session describe the warm-up transient rather than steady operation. A percentile over the whole session is displaced below the range the aircraft occupies once settled. \\
\hline
Drift / relative & Barometric height, derived height & Deviation from the metric's own trailing \SI{60}{\second} mean, with the threshold at the 99th percentile of the pooled Phase 1 deviation distribution & The absolute reading tracks atmospheric pressure rather than the aircraft, so a pooled percentile band would encode the prevailing weather. Deviation from the metric's own recent mean is comparable across sessions; the raw value is not. \\
\hline
Derived / variance & Attitude and acceleration instability & Rolling standard deviation over a \SI{5}{\second} window, with the threshold at the 99th percentile of the Phase 1 rolling-standard-deviation distribution & Instantaneous attitude values mean little on their own during normal manoeuvring. Dispersion is what separates controlled flight from loss of stability. \\
\hline
Liveness & Target reachability, metric staleness & Binary: \texttt{up == 0} or \texttt{absent()} sustained beyond two scrape intervals & Missing data is itself a fault condition, and cannot be expressed as a numeric bound. \\
\hline
\end{tabular}
\end{table}

Every rule also uses a Prometheus \texttt{for} condition, requiring the
threshold to remain active for 5 seconds before an alert is triggered. This
reduces false alarms caused by short-term fluctuations, which are common in
Wi-Fi SNR readings, but adds 5 seconds to the detection time. This additional
delay is considered in the latency targets presented in
Table~\ref{tab:eval-rq2}.

Time-of-flight distance is measured and reported but no threshold is derived
from it. Its upper tail arises from the aircraft passing over lower floor area in
one particular room, which is a property of the venue rather than of the
platform and would not transfer to another flight space.

The threshold values themselves are reported with the Study 2 results in
Section~\ref{sec:rq2}, since they are findings rather than design decisions.

\subsubsection{Data Analysis Plan}~\label{subsec:method-analysis}

Four measures are computed from the session logs. Each is defined here so that
any figure reported in Section~\ref{sec:rq2} can be traced back to the raw data
and recalculated.

\vspace{0.5em}
\noindent\textbf{Detection outcome.} Every fault run is classified into one of
three outcomes. A \textit{true positive} means the rule matching the induced
fault fired. A \textit{false negative} means no rule fired. A
\textit{misclassification} means a rule fired, but for the wrong fault class.
The nine runs are collected into a confusion matrix, which gives the detection
rate and precision.

\vspace{0.5em}
\noindent\textbf{Detection latency.} The gap between the fault appearing in the
telemetry and the alert firing, calculated as
$t_{\text{alert}}-t_{\text{onset}}$. Reported as a median, an interquartile
range and a maximum across all detected faults.

\vspace{0.5em}
\noindent\textbf{False positives.} An alert that fires while the drone is
behaving normally. This can happen in two places: during a baseline session, or
during a fault run before the fault has been induced. Both are counted. The
result is expressed as a rate, false alerts per \SI{10}{\minute} of normal
flight, rather than a raw count, so the figure stays comparable if the number of
sessions changes.

The rate is measured on normal flights held out from threshold derivation,
together with the pre-induction segment of each fault run. Measuring it on the
derivation sessions themselves would be circular, since a 99th-percentile
threshold is exceeded by approximately one percent of its own derivation data by
construction.

\vspace{0.5em}
\noindent\textbf{Metric coverage.} The proportion of monitored metrics that are
present and up to date in Prometheus for the full duration of a session, checked
with \texttt{up} and per-metric \texttt{absent()} queries. This confirms the
other three measures were computed from complete data.

\vspace{0.5em}
Only descriptive statistics are used. Five baseline sessions and nine fault runs
are too few to support significance testing, so results are reported as medians
with ranges, accompanied by the individual per-run figures. This lets the reader
judge the variability directly rather than relying on a summary statistic the
sample cannot justify. The small sample is a consequence of the experimental
design, and its effect on what can be claimed is addressed in
Section~\ref{subsec:method-threats}.

\begin{table}[htbp]
\centering
\caption{Study 2 Evaluation Criteria (RQ2)}
\label{tab:eval-rq2}
\small
\begin{tabular}{|p{2.8cm}|p{5.4cm}|p{2.6cm}|p{3.0cm}|}
\hline
\textbf{Criterion} & \textbf{Operational definition} & \textbf{Target} & \textbf{Measured by} \\
\hline
Threshold traceability & Every alerting rule threshold derived from Phase 1 baseline data by a stated rule, with no assumed values & 100\% of rules & Threshold table vs.\ baseline corpus \\
\hline
Detection rate & Induced faults for which the correct alerting rule fired & $\geq$ 8 of 9 runs & Confusion matrix \\
\hline
Detection latency & $t_{\text{alert}}-t_{\text{onset}}$ & Median $\leq$ \SI{5}{\second}; max $\leq$ \SI{10}{\second} & Alert log vs.\ raw telemetry log \\
\hline
False positive rate & False alerts per \SI{10}{\minute} of normal flight, pooled across baseline sessions & $\leq$ 1 per \SI{10}{\minute} & Alert log vs.\ baseline sessions \\
\hline
Metric coverage & All monitored metrics present and non-stale for the full duration of every session & $\geq$ 99\% of samples & \texttt{up}, \texttt{absent()} \\
\hline
Dashboard completeness & All monitored metrics, current alert state and fault annotations visible on one screen & All metrics and alert state & Dashboard inspection \\
\hline
\end{tabular}
\end{table}

The measures above define what is calculated. Table~\ref{tab:eval-rq2} states
the target each must meet for the study to be judged successful, fixed before
any data was collected.

All criteria in both studies were fixed before the corresponding data was
collected, and are expressed as numeric, observable targets. This is a
requirement of DSR, where artefact evaluation is subjective unless criteria are
defined in advance \citep{hevner2004}.

\subsection{Timeline and Milestones}~\label{subsec:method-timeline}

Figure~\ref{fig:gantt} shows the schedule for the remainder
of the trimester and Table~\ref{tab:milestones} states the milestones against
which progress is tracked.

% ---------------------------------------------------------------------------
% Move your existing Gantt chart (old Figure 4, RQ2 Milestone Plan) here.
% Shade completed bars and mark the current week.
% ---------------------------------------------------------------------------

\begin{table}[htbp]
\centering
\caption{Research Milestones}
\label{tab:milestones}
\small
\begin{tabular}{|p{0.6cm}|p{4cm}|p{5cm}|p{2.2cm}|p{1.7cm}|}
\hline
\textbf{\#} & \textbf{Milestone} & \textbf{Deliverable / evidence} & \textbf{Target} & \textbf{Status} \\
\hline
M1 & Pipeline validated on labelled data & Study 1 results, Section~\ref{sec:rq1} & T1 & Complete \\
\hline
M2 & Live Tello telemetry reaching Prometheus & Working exporter; target health page & 19 Jul 2026 & Complete \\
\hline
M3 & Baseline data collection complete & 5 archived session logs & Sun 30 Aug 2026 & Complete \\
\hline
M4 & Thresholds derived and rules encoded & Threshold table; \texttt{rules.yml} & Sun 30 Aug 2026 & Complete \\
\hline
M5 & Fault injection complete & 9 annotated run logs & Sun 7 Sep 2026 & \\
\hline
M6 & Analysis complete & Confusion matrix; latency and FP results & Fri 11 Sep 2026 & \\
\hline
M7 & Full draft to supervisor & Complete thesis draft & Mon 14 Sep 2026 & \\
\hline
M8 & Final submission & Final thesis and presentation & Mon 21 Sep 2026 & \\
\hline
\end{tabular}
\end{table}

\begin{table}[htbp]
\centering
\caption{Risks and Mitigations}
\label{tab:risks}
\footnotesize
\begin{tabular}{|p{3.2cm}|p{1.3cm}|p{6.4cm}|p{2.6cm}|}
\hline
\textbf{Risk} & \textbf{L / I} & \textbf{Mitigation} & \textbf{Status} \\
\hline
Threshold miscalibration causing false positives & High / High & Thresholds derived from Phase 1 baseline data by the rules in Table~\ref{tab:threshold-rules}, never assumed; \texttt{for:} hold durations suppress transients; thresholds pooled across all five baseline sessions rather than fitted to one; derived values validated against normal flights held out from the derivation & \textbf{Realised} in preliminary testing (\SI{60}{\celsius} assumption) and confirmed against baseline data \\
\hline
Metric assigned to the wrong behavioural class & Med / High & Metric class assignments in Table~\ref{tab:threshold-rules} tested by applying the derived threshold to held-out normal flights, since a misclassification is not visible from the derivation rule alone & \textbf{Realised} for barometric height and for onboard temperature; both reclassified \\
\hline
Tello hardware failure or unavailability & Med / High & Drone sourced from the Deakin equipment pool early in the trimester; spare batteries maintained; fallback to extended dataset-based experimentation preserves RQ2.1 and RQ2.2 & Monitored \\
\hline
Wi-Fi interference degrading telemetry & Med / Med & Flights conducted on the Tello's own access point indoors; multiple sessions per phase; outlier sessions reported rather than silently discarded & Controlled \\
\hline
Scrape interval too coarse for real-time detection & Low / Med & Scrape interval reduced from \SI{15}{\second} to \SI{1}{\second} for Study 2; storage overhead measured and reported & Mitigated \\
\hline
Insufficient fault range achievable indoors & Med / Med & Wi-Fi degradation supplemented by deliberate line-of-sight occlusion; outdoor contingency sessions available under CASA compliance & Monitored \\
\hline
Small sample limiting statistical claims & High / Low & Descriptive statistics only, with per-run figures reported; limitation stated explicitly & Accepted \\
\hline
Time constraints in a single trimester & Med / High & Two-week buffer between M6 and M7; Study 2 scoped to three fault classes & Monitored \\
\hline
\end{tabular}
\end{table}

Should baseline or fault collection be prevented entirely, the contingency in
Table~\ref{tab:risks} applies: analysis reverts to an extended dataset-based
experiment, which would narrow the contribution to RQ2.1 and RQ2.2 without
invalidating the design.

\subsection{Ethical Considerations}~\label{subsec:method-ethics}

\subsubsection{Human Participants and Ethics Determination}

This research involves no human participants, no personal data and no
observation of people. All data is machine-generated telemetry from a consumer
drone, so the project does not constitute human research under the National
Statement on Ethical Conduct in Human Research and requires no approval from
either a Faculty Human Ethics Advisory Group (low-risk review) or the Deakin
University Human Research Ethics Committee (DUHREC).

\subsubsection{Flight Safety and Regulatory Compliance}

Operating a powered aerial vehicle carries physical risk to people and property,
which is managed as follows.

\noindent\textbf{Indoor flights (primary).} All flights contributing data to
this study were conducted indoors at a private residence. Two exploratory
sessions were flown outdoors on 28 August 2026 under the CASA conditions
described below; both were excluded from analysis because baseline and
fault-injection runs must share a single environment, and both are retained in
the repository as rejected sessions. Civil Aviation Safety Authority (CASA)
rules governing outdoor airspace do not apply indoors, so for the sessions used
here the controlling considerations are venue and bystander safety. Permission to fly is obtained from the space's
custodian before each session. Flights are conducted in a cleared area with no
uninvolved persons present, propeller guards fitted at all times, and altitude
limited by ceiling clearance. The Tello's automatic low-battery landing failsafe
stays enabled in every session, including battery depletion runs, where it is the
object of measurement rather than a constraint to disable.

\noindent\textbf{Outdoor flights (contingency).} Should outdoor sessions be
needed to obtain a wider Wi-Fi degradation range, they will comply with CASA
recreational rules for very small remotely piloted aircraft (the Tello's take-off
weight is approximately \SI{80}{\gram}): below \SI{120}{\metre} AGL, in daylight,
within visual line of sight, clear of populous areas and emergency operations,
and outside the exclusion radius of controlled aerodromes. Current CASA guidance
will be checked immediately before any outdoor session rather than relied on from
this document, as the rules are subject to amendment. Outdoor sessions will not
be flown in rain or in wind exceeding the manufacturer's stated limit.

\subsubsection{Data Management and Privacy}

No personally identifiable information is collected at any stage. The Tello
carries no GNSS receiver, so no location data is recorded, and although the
airframe has a camera, the video stream is not enabled, captured or stored in any
session. Only numeric telemetry is collected.

Telemetry, session logs and configuration files are stored locally on
password-protected, full-disk-encrypted hardware, with backups to a private
version-controlled repository accessible only to the researcher and supervisor
during the project. Data is retained for the minimum period required by
university policy following submission. On submission the telemetry and
configuration artefacts are released publicly as described in
Section~\ref{subsec:method-summary}; because the data contains no personal or
sensitive information, no de-identification is required before release.

\subsubsection{Research Integrity and Use of Generative AI}

All results are reported as observed, including failures, missed detections and
false positives. The \SI{60}{\celsius} false positive encountered during
preliminary testing is reported in full instead of discretely corrected, because
it is evidence for the central argument of RQ2. Negative and partial results are
reported with the same prominence as successful ones.

Generative AI tools were used in a defined and limited capacity, consistent with
Deakin University's requirements on acknowledging AI assistance. Claude
(Anthropic) was used to assist in producing the software instruments for both
studies: the Study 1 replay exporter, and the Study 2 session harness, session
validator and threshold derivation script. It was also used for code-level
debugging and document structuring. The research design, research questions,
choice of methods, metric selection, threshold derivation rules, experimental
protocol and all interpretation of results are the author's own work.

Two distinctions are stated explicitly. First, the derivation rules in
Table~\ref{tab:threshold-rules} were specified by the author; the derivation
script implements those rules and computes their outputs from the measured data,
and no threshold value was proposed or selected by an AI tool. Second, no AI tool
generated, altered or fabricated any telemetry. All data derives from physical
flights, and the raw session logs are published so that every reported figure can
be recomputed independently. All AI-assisted code was reviewed, executed and
validated by the author against real flights before use.

\subsection{Threats to Results}~\label{subsec:method-threats}

Threats are organised by type, following standard practice in experimental
software engineering \citep{wohlin2012}. Risks to the project schedule are
handled separately in Table~\ref{tab:risks}.

\noindent\textbf{Internal validity.} Battery ageing, cell condition and ambient
temperature vary between sessions and affect both depletion rate and thermal
readings, so battery identifier and ambient temperature are logged per session
and any outlier can be attributed rather than dismissed. More significantly,
thresholds are derived from baseline flights on the same airframe later used for
fault runs, which risks fitting the thresholds to that specific unit. Pooling
across five baseline sessions rather than one reduces this, but it cannot be
eliminated with a single available drone, and it constrains the strength of the
RQ2.1 claim.

\noindent\textbf{Ambient temperature.} The temperature threshold is expressed as
an absolute value and is therefore sensitive to the conditions under which the
baseline was collected. All five baseline sessions were flown at an ambient
temperature of \SI{21}{\celsius} and their settled means fall between
\SI{61.4}{} and \SI{62.8}{\celsius}. A held-out normal flight conducted in a
room estimated to be \SI{1}{} to \SI{2}{\celsius} warmer settled at
\SI{64.0}{\celsius}, above every baseline session, while its battery depletion
rate and mean acceleration magnitude remained within the baseline range. A
baseline collected at a single ambient temperature cannot quantify that
relationship, so the derived value transfers only to comparable conditions.

\noindent\textbf{Positioning and ambient lighting.} The aircraft holds position
using a downward-facing camera and a time-of-flight sensor, and therefore depends
on ambient lighting. A session attempted in a dimly lit room produced
time-of-flight readings ranging between \SI{21}{} and \SI{133}{\centi\metre}
against a stable \SI{77}{} to \SI{78}{\centi\metre} in adequate light,
uncommanded lateral drift, refusal of relative-position commands, and a battery
depletion rate approximately \num{2.5} times the baseline. Lighting is therefore
controlled alongside venue and ambient temperature, and the affected session was
excluded.

\noindent\textbf{External validity.} The study uses one drone model in a
predominantly indoor environment with a small number of sessions, so the specific
threshold values are not transferable to other platforms. The Tello's telemetry
is also narrower than that of the platforms in much of the reviewed literature,
providing no GNSS data and only coarse Wi-Fi quality information, which limits
how directly the findings compare to GPS-equipped research platforms. What is
claimed as transferable is the derivation \textit{procedure} and the pipeline
architecture, not the numbers, and that distinction is maintained throughout the
reporting of results.

\noindent\textbf{Construct validity.} A composite health-status metric compresses
a multi-dimensional condition into a single value, and whether it represents
"health" in a way an operator would accept is a modelling assumption rather than
a measured fact. Detection latency is bounded below by the scrape interval and
the \texttt{for:} hold duration, so the figures reported measure the latency of
\textit{this configuration}, not an intrinsic property of threshold-based
detection. Both bounds are stated alongside every latency result. Defining fault
onset from the first raw threshold crossing partially couples the onset
definition to the threshold under test; the operator induction marker is kept as
an independent check against this.

\noindent\textbf{Reliability.} All session logs, configuration files, alert
histories and PromQL queries are archived and published, so every reported figure
can be independently recomputed from the raw data. Tool versions are pinned.
Where a session is excluded from analysis, the exclusion and its reason are
reported rather than the session being silently dropped.

\subsection{Summary}~\label{subsec:method-summary}

The methodology is one experimental design executed as two studies against a
single artefact. Study 1 established that the pipeline can observe health state,
using a pre-labelled dataset so that the pipeline itself was the only thing under
test. Study 2 removes the labels and requires the same pipeline to determine
health state from raw telemetry on physical hardware, with every threshold
derived from the drone's own measured baseline rather than assumed. Experimental
research supplies the controlled conditions that make detection latency and false
positive rate measurable; design science research supplies the structure for
evaluating the pipeline as a contribution in its own right. Evaluation criteria
for both studies were fixed before the corresponding data was collected, and the
limits of what the results can support are stated in
Section~\ref{subsec:method-threats} rather than left implicit.

\noindent\textbf{Sustainability and reuse.} All research artefacts are published
in a public repository on submission
(\url{https://github.com/Dmatt547/RAS-monitoring-thesis.git}): both exporters,
\texttt{prometheus.yml}, \texttt{rules.yml} with the derived thresholds, the
Grafana dashboard as exported JSON, the full baseline and fault-run telemetry
logs, the PromQL queries used in analysis, and a README documenting the
reproduction procedure end to end. The Study 1 dataset is not redistributed, as
it is licensed CC BY-SA 4.0 \citep{ekanayakadevlk2024}; it is cited and linked
instead, with the preprocessing script provided so the replay can be
reconstructed. Tool versions are pinned in Table~\ref{tab:resources} and the
deployment is containerised, so the pipeline can be rebuilt without manual
environment configuration.

Every component of the stack is open-source, with no paid licence, subscription
or cloud service in the dependency chain. That is a deliberate design property,
since the gap addressed is the absence of monitoring approaches that survive
without cloud infrastructure, and it means the artefact keeps working after the
project ends and after any university account access ends. The exporter
presents a platform-agnostic interface, so any system reporting numeric telemetry
can replace the Tello without altering the Prometheus, alerting or dashboard
layers. The derivation rules in Table~\ref{tab:threshold-rules} are stated in
terms of metric behaviour rather than drone-specific fields. Prometheus'
multi-target scrape model also allows extension from one drone to a fleet without
architectural change. These extensions are discussed in
Section~\ref{subsec:future-work}.
